\chapter{Model Merging and Recombination}
\minitoc

\begin{chapteroverview}
  \begin{itemize}
    \item Master task vector arithmetic and SLERP weight interpolation.
    \item Implement TIES-Merging and DARE for combining fine-tuned experts.
    \item Evaluate the "FrankenMoE" strategy for creating heterogeneous ensembles.
    \item Analyze the 2025 shift from training singular models to evolving model families.
  \end{itemize}
\end{chapteroverview}

Model merging combines the weights of multiple fine-tuned models without additional training, enabling capability combination at near-zero compute cost.

\section{Merging Methods}

\begin{itemize}
  \item \textbf{Linear interpolation:} $\theta = (1-\alpha)\theta_A + \alpha\theta_B$. Simple but creates interference between dissimilar capabilities.
  \item \textbf{SLERP (Spherical Linear Interpolation):} Interpolates along the geodesic of weight space. Better than linear for preserving each model's characteristics.
  \item \textbf{TIES-Merging \parencite{yadav2023ties}:} Trim redundant deltas $\rightarrow$ Elect sign direction $\rightarrow$ Disjoint merge. Handles sign conflicts between different fine-tuned models.
  \item \textbf{DARE \parencite{yu2023language} (Drop And REscale):} Randomly drop task-vector parameters and rescale survivors. Reduces interference between merged models by eliminating redundant updates.
  \item \textbf{Model Breadcrumbs:} Remove outlier parameters before merging. Outliers cause interference; removing them improves merge quality.
\end{itemize}

\section{Task Vectors}

Full formulations in Appendix~\ref{app:merging}: \hyperref[form:task_vec]{task arithmetic}, \hyperref[form:slerp]{SLERP}, \hyperref[form:ties]{TIES}, \hyperref[form:dare]{DARE}; code: \hyperref[lst:task_vec]{Listing~\ref*{lst:task_vec}}, \hyperref[lst:ties]{Listing~\ref*{lst:ties}}.

Task vector \parencite{ilharco2022editing} $\tau = \theta_\text{fine-tuned} - \theta_\text{base}$. Arithmetic on task vectors enables:
\begin{itemize}
  \item \textbf{Capability addition:} $\theta_\text{base} + \tau_A + \tau_B$ combines coding and math capabilities.
  \item \textbf{Capability subtraction:} $\theta_\text{base} - \tau_\text{toxic}$ removes unwanted behaviors.
  \item \textbf{Analogy:} $\theta_\text{base} + \tau_\text{French} - \tau_\text{English}$ adapts a model to French.
\end{itemize}

\begin{keyinsight}[When Merging Works Best]
Merging works best when models share the same base and have been fine-tuned on complementary (not conflicting) tasks. Models fine-tuned on the same task but different data benefit more from ensemble methods than weight merging.
\end{keyinsight}

\section{Practical Merging Workflow}

\begin{enumerate}
  \item Fine-tune $n$ specialized models from the same base.
  \item Compute task vectors $\tau_i = \theta_i - \theta_\text{base}$ for each.
  \item Apply TIES or DARE to reduce interference.
  \item Evaluate merged model on target benchmarks; tune merge coefficients $\alpha_i$.
  \item Distribute as a single merged model.
\end{enumerate}

Tools: \texttt{mergekit} \parencite{ilharco2022editing} (open source), Hugging Face \texttt{transformers} weight manipulation.
